\section{Introduzione all'interpretazione astratta}

\subsection{Semantica astratta sui CFG}

Mentre nei capitoli precedenti è stata introdotta la semantica
collettrice sul programma sorgente, per l'analisi statica è più conveniente
ragionare in termini di CFG.
In questa rappresentazione, la semantica dei costrutti di controllo come
\texttt{while} e \texttt{if} è codificata direttamente nella struttura del
grafo stesso.

Dato il seguente frammento di codice:
\begin{verbatim}
1 x = 0;
2 while ( x < 100 )
3   if ( x < 50 )
4       x = x + 2;
5   else
6       x = x + 10;
\end{verbatim}

È possibile costruire il seguente CFG (\cref{fig:th_08_01}):

\begin{figure}[H]
    \centering
    \includegraphics[width=0.6\textwidth]{images/th_08/01.png}
    \caption{L'immagine illustra il CFG corrispondente al frammento di codice
    mostrato.}
    \label{fig:th_08_01}
\end{figure}
Lo scopo è di definire gli stati di ingresso e uscita per ogni nodo del CFG
utilizzando la semantica astratta.

\begin{definizione}{Entry state}{entry_state}

Per ogni punto del programma (nodo del CFG) $x_i$, lo stato di ingresso
(\textit{entry}) è definito come il \textit{least upper bound}
(\Cref{def:upper_bound}) degli stati di uscita dei blocchi che puntano ad
esso (i predecessori). Per il blocco iniziale, è lo stato di ingresso dato.

\end{definizione}

\begin{definizione}{Exit state}{exit_state}

Per ogni punto del programma (nodo del CFG) $x_i$, lo stato di uscita
(\textit{exit}) è il risultato dell'applicazione della semantica (concreta o
astratta) dello statement contenuto nel blocco, applicata allo stato di
ingresso.

\end{definizione}

Poiché la logica di controllo è gestita dagli archi del grafo, è sufficiente
esprimere la semantica solo per gli assegnamenti e le espressioni booleane,
che fungono da ``guardie'' per gli archi.


\begin{esempio}{Calcolo degli stati di ingresso}{cfg_entry}
Prendendo in considerazione il CFG in \Cref{fig:th_08_01}, è possibile costruire
gli stati di ingresso per ogni blocco come segue:
\[
\begin{aligned}
    x_1 & = I \\
    x_2 & = [\![ x = 0 ]\!]x_1 \cup [\![ x = x + 2 ]\!]x_4 \cup [\![ x = x + 10 ]\!]x_5 \\
    x_3 & = [\![ x < 100 ]\!]x_2 \\
    x_4 & = [\![ x < 50 ]\!]x_3 \\
    x_5 & = [\![ x \geq 50 ]\!]x_3 \\
    x_6 & = [\![ x \geq 100 ]\!]x_2 \\
\end{aligned}
\]
\end{esempio}

\subsection{Sistemi di Equazioni e punto Fisso}

Dato un CFG con $m$ nodi si vorrebbe procedere alla risoluzione di un sistema
di equazioni.

\begin{definizione}{Sistema di Equazioni Semantiche}{semantic_equations}
Il sistema è definito come:
\[ F = \{ x_i = f_i(x_1, \dots, x_m) \mid i = 1, \dots, m \} \]
dove ogni $x_i$ rappresenta lo stato (o la proprietà) associata al nodo $i$.
\end{definizione}

L'obiettivo è calcolare la soluzione minima del sistema $F$ (\Cref{def:semantic_equations}),
che corrisponde al limite di un'iterazione di Kleene sul reticolo del dominio:
$\langle \wp(\mathbb{Z}), \subseteq, \emptyset, \mathbb{Z}, \cup, \cap \rangle$.

\begin{nota}{Punto fisso sul dominio concreto}{concrete_fixed_point}
Calcolare il least fixpoint (lfp) sul dominio concreto ($\wp(\mathbb{Z})$) è
spesso indecidibile o non computabile in tempo finito.
\end{nota}

\subsection{Approssimazione del punto fisso}

Dato che molto spesso non è possibile calcolare il punto fisso sul dominio concreto,
si introduce un dominio astratto $A$ che approssima il dominio concreto $C$.

L'obiettivo è calcolare un'approssimazione corretta $\alpha(\text{lfp}(F))$
senza dover calcolare effettivamente $\text{lfp}(F)$.

Per farlo, è necessario definire un sistema di equazioni astratto $F^\#$ e
calcolarne il suo punto fisso.

\subsubsection{Soundness}

Affinché l'analisi sia affidabile, la funzione astratta $f^\#$ deve
essere una sound abstraction della funzione concreta $f$.


\begin{definizione}{Soundness della funzione astratta}{abstract_soundness}

Siano $C$ e $A$ due domini legati da una connessione di Galois $(\alpha, \gamma)$
e siano $f: C \to C$ una funzione concreta e $f^\#: A \to A$ una funzione astratta.
Allora, $f^\#$ è una sound abstraction di $f$
\[ \forall c \in C . \alpha(f(c)) \sqsubseteq_A f^\#(\alpha(c)) \]
Oppure equivalentemente:
\[ \forall a \in A . f(\gamma(a)) \sqsubseteq_C \gamma(f^\#(a)) \]

\begin{figure}[H]
    \centering
    \includegraphics[width=0.8\textwidth]{images/th_08/02.png}
    \caption{L'immagine illustra la soundness della funzione astratta.}
    \label{fig:th_08_02}
\end{figure}

\end{definizione}

\begin{nota}{Best correct approximation (BCA)}{bca}
Tra tutte le possibili funzioni astratte corrette, esiste la ``migliore''
(la più precisa), definita come:
\[ f^{bca} = \alpha \circ f \circ \gamma \]
Tuttavia, spesso la BCA non è computabile perché richiede l'applicazione di
$f$ e $\alpha$, operazioni che potrebbero essere troppo complesse.
\end{nota}

\subsection{Teoremi di trasferimento del punto fisso}

Una volta definito il sistema astratto $F^\#$, esistono due teoremi che
garantiscono che il suo punto fisso $lfp(F^\#)$ approssimi correttamente il punto
fisso concreto $lfp(F)$.

\begin{nota}{Approssimazione e esattezza}{approximation_exactness}

$lfp(F^\#)$ è una approssimazione soudn di $lfp(F)$ se:
\begin{itemize}
    \item Esattezza: $\alpha(\text{lfp}(F)) = \text{lfp}(F^\#)$
    \item Approssimazione: $\alpha(\text{lfp}(F)) \sqsubseteq \text{lfp}(F^\#)$
\end{itemize}

\end{nota}

\begin{teorema}{Approssimazione di Kleene}{kleene_approx}

Se $f: C \to C$ è continua su un reticolo completo concreto $\langle C, \sqsubseteq_C, \bot_C, \top_C, \sqcup_C, \sqcap_C \rangle$
e $f^\#: A \to A$ è un'astrazione sound di $f$ su un dominio astratto
$\langle A, \sqsubseteq_A, \bot_A, \top_A, \sqcup_A, \sqcap_A \rangle$, e se
la sequenza $\{f^\#_i(\bot_A) \mid i \in \mathbb{N} \}$ ha un limite $x \in A$,
allora è una approssimazione sound di $\text{lfp}(f) \sqsubseteq_C \gamma(x)$

In particolare, ad ogni passo $i$, l'iterazione astratta $f^\#_i(\bot_A)$
è una approssimazione sound di $f_i(\bot_C)$.

\end{teorema}


\begin{nota}{Approssimazione di Kleene}{}

Questo teorema garantisce che l'iterazione di Kleene sul dominio astratto
fornisce una approssimazione sound del punto fisso concreto.
Tuttavia, non garantisce che l'approssimazione sia la migliore possibile
(neanche alla convergenza).

\end{nota}

\begin{teorema}{Approssimazione di Tarski}{tarski_approx}

Dato un reticolo completo $\langle C, \sqsubseteq_C, \bot_C, \top_C, \sqcup_C, \sqcap_C \rangle$
e una funzione monotona $f: C \to C$, e una sua astrazione sound $f^\#: A \to A$
in un dominio astratto $\langle A, \sqsubseteq_A, \bot_A, \top_A, \sqcup_A, \sqcap_A \rangle$,
allora qualsiasi \textbf{post-fixpoint} $a$ di $f^\#$ (ovvero tale che $f^\#(a) \sqsubseteq_A a$)
è un'approssimazione sound del punto fisso concreto $lfp(f)$, ovvero:
\[ \text{lfp}(f) \sqsubseteq_C \gamma(a) \]

\end{teorema}
\begin{nota}{Approssimazione di Tarski}{}

Il teorema di approssimazione di Tarski garantisce che qualsiasi
post-fixpoint della funzione astratta è una approssimazione sound del
punto fisso concreto.

\end{nota}
